% =====================================================================
\section{WiFiManager: настройка сети без перепрошивки}\label{les:wifimanager}
% =====================================================================
\lessonintro
  {дать пользователю ограниченный по времени способ настроить Wi-Fi и сохранить локальную работу}
  {уроки 8 (STA/AP), 21 (автоматы), 38 (первичная настройка), 39 (ESPUI)}
  {портал по удержанию кнопки и подключение к выбранной пользователем сети}
  {WiFiManager 2.0.17, кнопка GPIO5--GND, LED через резистор на GPIO4, сеть 2,4 ГГц}

\subsection{Панель управления и портал настройки}

ESPUI отвечает на вопрос «как управлять устройством». WiFiManager --- «как сообщить ему
имя и пароль сети, не редактируя исходники». Это разные этапы жизненного цикла.
Provisioning нужен при первом запуске, переезде в другую сеть или осмысленном запросе владельца,
а обычная панель нужна во время работы.

Портал создаёт временную точку доступа и страницу выбора сети. Captive portal пытается
направить подключённый браузер на страницу настройки, но автоматическое окно зависит от
телефона и ОС. Всегда оставляйте способ открыть IP вручную.
У ESP32-S3 поддерживается Wi-Fi 2,4 ГГц; наличие сети только 5 ГГц не исправляется библиотекой.

\subsection{Когда открывать доступ}

Не стирайте учётные данные при каждой потере Wi-Fi и не оставляйте открытую AP навсегда.
В нашем опыте плата пробует ранее сохранённую сеть. Для настройки пользователь удерживает
физическую кнопку 2 секунды; открывается окно с целевым временем 120 секунд.
После отпускания кнопки можно запросить новую попытку. Пароль \code{setup-40-demo} учебный:
на собственных платах задавайте индивидуальный и сообщайте владельцу отдельно.

Пароль AP и пароль выбранной Wi-Fi-сети --- разные секреты. Временная сеть
\code{S3-Setup} служит для настройки, а не для выдачи доступа в школьную административную сеть.
Учётные данные выбранного роутера сохраняет стек Wi-Fi по используемому библиотекой пути.
Новые поля приложения автоматически такими данными не становятся: для них отдельно
проектируют валидацию, версию формата и запись в Preferences.

\subsection{Неблокирующий режим и явное окно времени}

Установите WiFiManager через менеджер библиотек. В главе проверена версия 2.0.17.
Объект объявлен глобально: при неблокирующем режиме он должен существовать и после
завершения \code{setup()}. Функция \code{process()} обслуживается из основного цикла.
Закрытие окна реализуем своим счётчиком \code{millis()}, а не полагаемся на действие
настроек тайм-аута из другого режима библиотеки. Закрытие проверяется в цикле: сетевой вызов
может задержать его, поэтому это не жёсткая гарантия времени.

\codecap{Портал по кнопке, собственный тайм-аут и независимое мигание}
\codeinput[\footnotesize]{l40_wifimanager}

Подключите LED к GPIO4 через резистор, кнопку --- к GPIO5 и GND.
После старта удержите кнопку, подключитесь к \code{S3-Setup} и откройте адрес из лога.
Выберите тестовую сеть 2,4 ГГц и введите её пароль. Затем вернитесь на эту сеть
и проверьте выданный устройству IP. Перезапустите плату: ранее сохранённые данные
должны позволить подключение без повторного ввода.

Имя «неблокирующий» не обещает жёсткого времени каждого вызова. Сканирование,
подключение и обработка HTTP имеют собственную стоимость; в примере задан предел попытки
подключения. Мигание показывает, что программа в целом продолжает работать, но не служит
доказательством точного периода во время всех сетевых операций.
Для аппаратно критичного действия нужны соответствующие таймер, задача и независимая защита.

\subsection{Как совместить с ESPUI}

Оба средства по умолчанию хотят обслуживать HTTP на порту 80. Нельзя просто объединить
два скетча и дважды открыть сервер на одном адресе и порту.
Для первого проекта используйте последовательные этапы:

\begin{enumerate}
  \item устройство принимает запрос настройки и запускает портал WiFiManager;
  \item после успешной настройки портал закрывается и его сетевые ресурсы освобождаются;
  \item выполняется перезапуск либо явный переход режима, затем запускается ESPUI;
  \item повторная настройка снова переводит устройство в отдельный режим.
\end{enumerate}

Состояние режима храните отдельно от переключателя LED. Перезапуск после сохранения ---
простой учебный способ разделить серверы; запускайте его после подтверждения сохранения,
а не в середине HTTP-обработчика. При потере сети в рабочем режиме устройство сохраняет
локальное управление и пробует восстановить связь, а не немедленно возвращается к настройке.

Если нужен одновременный доступ, придётся осознанно спроектировать владение серверами,
портами, DNS и ресурсами: это отдельная архитектура. Captive portal не заменяет авторизацию
рабочей панели и не добавляет HTTPS. Сначала проверьте каждый пример отдельно.

\subsection{План проверки владельцем устройства}

Проверьте первое включение без настроек; неверный пароль; закрытие 120-секундного окна;
выключенный роутер; повторный запрос по кнопке; перезапуск после успешной настройки.
Зафиксируйте ожидаемое безопасное состояние нагрузки и максимальное допустимое время ожидания.
Не выводите пароли в Serial, веб-журнал или публичный репозиторий.
Стирание всех настроек --- отдельная команда с подтверждением, а не штатное действие при старте.

\begin{trouble}
\begin{itemize}[leftmargin=*]
  \item Не открылось автоматическое окно --- введите IP вручную; проверьте, что телефон остался в AP.
  \item Нет нужной сети --- убедитесь в диапазоне 2,4 ГГц и уровне сигнала.
  \item Программа перестала обслуживать портал --- объект уничтожен или не вызывается process().
  \item ESPUI не запускается --- проверьте, остановлен ли сервер настройки и освобождён ли порт.
\end{itemize}
\end{trouble}

\begin{summary}
WiFiManager упрощает выдачу параметров сети. Приложение определяет, кто и когда может
открыть портал, сколько он работает и что происходит при ошибке.
Настройка, рабочая панель и восстановление связи --- отдельные состояния.
\end{summary}

\begin{selfcheck}
\begin{enumerate}
  \item Чем пароль временной AP отличается от пароля роутера?
  \item Почему объект WiFiManager в этом примере не локален для setup()?
  \item Почему два готовых веб-скетча нельзя просто запустить одновременно на порту 80?
\end{enumerate}
\end{selfcheck}

\needspace{20\baselineskip}
\begin{tasks}
\begin{enumerate}
  \item \lv{1} Настройте подключение к тестовой сети через портал и найдите IP в логе.
  \item \lv{1} Перезапустите плату и подтвердите подключение без повторного ввода.
  \item \lv{2} Проверьте неверный пароль и закрытие окна через 120 секунд; запишите фактическое время.
  \item \lv{2} Добавьте понятный индикатор режимов «настройка», «подключение», «работа» без бесконечного ожидания.
  \item \lv{3} Добавьте проверяемое поле имени устройства и сохраняйте его отдельно от учётных данных Wi-Fi.
  \item \lv{4} Объедините настройку и ESPUI в автомат с отдельными этапами, проверьте повторную настройку, потерю сети и восстановление без конфликта серверов.
\end{enumerate}
\end{tasks}

\paragraph{Документация для продолжения}
\href{https://github.com/tzapu/WiFiManager}{WiFiManager: API и примеры},
\href{https://github.com/tzapu/WiFiManager/tree/master/examples/NonBlocking}{неблокирующие сценарии}.
